% ── Chapter 3: Methodology ────────────────────────────────────
\chapter{Methodology}

\section{System Architecture Overview}

The proposed system is a full-stack web application following a client-server architecture
designed around the separation of concerns between the user interface, the AI inference
pipeline, and the data persistence layer. This architectural decomposition enables independent
development and testing of each component and allows the computationally intensive AI backend
to be scaled or upgraded without changes to the frontend user experience.

The frontend is developed in React, a component-based JavaScript framework built on HTML
\cite{sharma2018html,mozilla2025html}, and provides users with an image upload interface,
a live camera capture option, and an interactive results visualization dashboard. The backend
is implemented in Python \cite{siva2023python} using FastAPI \cite{tioman2024fastapi}, a high-performance asynchronous
web framework that exposes a RESTful API consumed by the React frontend.
 Upon
receiving an image via the API, the backend executes the full AI inference pipeline and returns
a structured JSON response containing all detection results, diameter measurements, OCR labels,
and susceptibility classifications. A relational database implemented in SQLite stores analysis
history, user session records, and the CLSI 2025 and EUCAST 2023 breakpoint tables, enabling
both historical audit of past analyses and rapid in-process lookup of breakpoint thresholds.

The overall image analysis pipeline consists of five sequential stages, each of which is
described in detail in the sections that follow:
\begin{enumerate}
  \item \textbf{Image acquisition and preprocessing:} The input image is validated, resized to
        the model's required input resolution, and preprocessed according to the requirements
        of the active detection model.
  \item \textbf{Antibiotic disk detection and OCR labeling:} The detection model localizes
        individual antibiotic disks; the OCR pipeline reads the alphanumeric label from each
        disk crop to identify the antibiotic.
  \item \textbf{ZOI detection and diameter measurement:} The detection model localizes each
        zone of inhibition as a bounding box; the bounding box dimensions are converted to a
        diameter estimate in pixels.
  \item \textbf{Pixel-to-millimeter calibration:} The pixel-scale diameter measurements are
        converted to physical millimeters using the known 6\,mm diameter of the antibiotic
        disk as an in-image reference object.
  \item \textbf{Breakpoint lookup and susceptibility classification:} The measured ZOI diameter
        in millimeters is compared against the CLSI or EUCAST breakpoint for the identified
        antibiotic and the user-specified organism, and the organism is classified as
        Susceptible (S), Intermediate (I), or Resistant (R).
\end{enumerate}

\begin{figure}[H]
  \centering
  \safeimg[width=\textwidth]{figures/pipeline_flowchart.png}{End-to-end analysis pipeline}
  \caption{End-to-end image analysis pipeline of the automated AST system, from image input
    through detection, OCR, calibration, and breakpoint-based susceptibility classification.}
  \label{fig:pipeline}
\end{figure}

\begin{figure}[H]
  \centering
  \safeimg[width=0.92\textwidth]{figures/system_architecture.png}{System Architecture Diagram}
  \caption{Overall system architecture of the automated AST web application, showing the
    React frontend, FastAPI backend, AI inference pipeline, and SQLite database layer.}
  \label{fig:sys_arch}
\end{figure}

\section{Dataset Collection and Preparation}

Assembling a high-quality, diverse training dataset is essential for the generalization
performance of deep learning models. For the AST domain, dataset construction is particularly
challenging because real clinical plate images are often proprietary (belonging to hospitals
or reference laboratories), and the volume of publicly available annotated AST image data is
limited. The dataset strategy employed in this project therefore combines three complementary
sources: real laboratory images, augmented real images, and programmatically generated
synthetic images.

\subsection{Real Image Collection}

A total of 287 real AST plate photographs were collected from three sources. Open sharing of
biomedical research data through repositories such as DRYAD has been a longstanding principle
in scientific practice \cite{iom1996}. The first source was the publicly available DRYAD data
repository \cite{giske2024dryad}, which contributed 225 annotated disk diffusion plate images
captured in various laboratory environments. The second source was the
CDS (Clinical and Demographic Study) dataset, contributing 50 images. The third source was
the CRA (Chulabhorn Royal Academy) dataset, contributing 32 images. Together these three
sources yielded the 287 real images used in this project.

Images span multiple bacterial species, a range of antibiotic disk counts per plate
(from 1 to 10 disks), and a variety of zone overlap configurations that represent the
diversity of clinical scenarios the system must handle. The externally sourced subsets
introduce controlled variability in lighting, camera model, agar brand, and plate size that
improves the generalization of models trained on the combined dataset.

\subsection{Data Augmentation}

To expand the training set from 287 images while preserving the label correctness of the
original annotations, a set of geometric and photometric augmentation operations were applied
to the real images using the Roboflow platform. Augmentation transformations included
horizontal and vertical flipping, random rotation within $\pm$15$^\circ$, brightness variation
($\pm$25\%), contrast variation ($\pm$20\%), and additive Gaussian noise (standard deviation
up to 5\% of the maximum pixel intensity). Each augmented image was generated by randomly
sampling from this set of transformations, with bounding box coordinates automatically
recomputed to match the transformed image geometry. This augmentation strategy increased the
training dataset from 287 to 543 images, approximately doubling the number of unique training
examples available from real plate photography.

It is important to note that augmentation was applied only to the training partition of the
real images; the validation and test partitions were held out from augmentation to ensure
that evaluation metrics reflect model performance on unmodified, realistic inputs.
Augmentation was also deliberately restricted to transformations consistent with plausible
variation in real laboratory photography, avoiding extreme transformations (e.g., heavy
perspective warping or color inversion) that would produce unrealistic images detrimental to
model training.

\subsection{Synthetic Data Generation}

An additional 2,000 synthetic images were generated programmatically to further increase
dataset diversity and to address the relative scarcity of real images containing specific
zone overlap configurations and atypical zone morphologies. Synthetic plate images were
rendered using a procedural pipeline implemented in Python, using Pillow and OpenCV for
image composition. Each synthetic image was designed to visually approximate the appearance
of a real Mueller-Hinton agar plate, with a textured agar background generated by tiling and
blending sampled texture patches, circular antibiotic disk objects placed at defined positions,
and circular zone-of-inhibition regions rendered with Gaussian-blurred edges to mimic the
gradual growth boundary seen in real photographs.

The 2,000 synthetic images were divided into two categories of equal size. The first
category of 1,000 images used fixed antibiotic disk positions arranged according to the
standard CLSI-recommended plate layout for common 6-disk and 8-disk configurations,
simulating the regular, structured patterns seen in standardized laboratory workflows.
The second category of 1,000 images used fully randomized disk placement within the plate
boundary, simulating the less structured configurations common in research and teaching
laboratories. Across both categories, zone diameters were sampled from a realistic
distribution spanning the full clinical range (from approximately 10\,mm to 40\,mm),
with deliberate over-representation of overlapping zone scenarios---where the inhibition
zones of adjacent disks partially intersect---to ensure the models were exposed to this
challenging but clinically important case during training. Background texture, simulated
lighting gradients, and synthetic colony growth patterns were also varied to increase
visual diversity.

\subsection{Final Dataset Split}

\begin{table}[H]
  \centering
  \caption{Dataset composition and train/validation/test split used for model training and evaluation.}
  \label{tab:dataset_split}
  \begin{tabular}{lcccc}
    \toprule
    \textbf{Source} & \textbf{Total} & \textbf{Train} & \textbf{Validation} & \textbf{Test} \\
    \midrule
    DRYAD                              & 225   & ---   & ---   & ---   \\
    CDS                                & 50    & ---   & ---   & ---   \\
    CRA (Chulabhorn Royal Academy)     & 32    & ---   & ---   & ---   \\
    \textit{Real total (original)}     & 287   & ---   & ---   & ---   \\
    \textit{Real total (augmented)}    & 543   & 433   & 55    & 55    \\
    Synthetic (fixed position)         & 1,000 & \multicolumn{3}{c}{training only} \\
    Synthetic (unfixed position)       & 1,000 & \multicolumn{3}{c}{training only} \\
    \bottomrule
  \end{tabular}
\end{table}

The synthetic images were used exclusively for training, as they were generated to augment
model learning on diverse configurations rather than to evaluate generalization performance.
All evaluation metrics reported in Chapter~4 are computed on the real-image test set (55 images)
to ensure that reported performance figures reflect the model's behavior on authentic,
unmanipulated laboratory photographs.

\section{Annotation}

Ground-truth annotations were created in Pascal VOC XML format, specifying bounding box
coordinates for each zone of inhibition and each antibiotic disk visible in every image.
This annotation format encodes four values per object: the pixel coordinates of the
top-left corner and the bottom-right corner of the bounding box, together with a class
label (\texttt{antibiotic} for disks, \texttt{disk-zone} for inhibition zones).

Initial annotations were created by the project team using LabelImg, an open-source graphical
image annotation tool. To ensure clinical correctness, all bounding box placements for the
zone-of-inhibition class were subsequently reviewed and corrected by a trained microbiologist
who independently verified each annotation against the corresponding physical plate measurement
taken with calipers. This dual-review workflow was essential for producing a reliable ground
truth, since zone boundary definition requires domain expertise: the correct edge of an
inhibition zone is defined as the point of complete inhibition of confluent growth (not the
zone of micro-colony formation, which CLSI instructs technicians to disregard). Annotations
were then imported into Roboflow for dataset management, augmentation pipeline configuration,
and export in the formats required by each model's training framework (Pascal VOC XML for
Faster R-CNN and YOLO TXT format for YOLOv11n).

\section{Model Training}

\subsection{Faster R-CNN}

Faster R-CNN was implemented using PyTorch's \texttt{torchvision} library with a ResNet-50
backbone pretrained on the ImageNet large-scale visual recognition dataset \cite{he2016deep}.
Using a pretrained backbone is critical for performance on a medical imaging dataset of this
size: the ImageNet pretraining initializes the backbone with discriminative feature representations
for low-level image structure (edges, textures, gradients) that transfer effectively to the
AST domain, significantly accelerating convergence and improving final accuracy compared to
training from random initialization.

The detection head (both the RPN and the classification/regression head) was initialized
from scratch and trained from the start. Training used the Stochastic Gradient Descent (SGD)
optimizer with a learning rate of $0.005$, momentum of $0.9$, and weight decay of
$5 \times 10^{-4}$ (i.e., \texttt{0.0005}) to regularize the model and prevent overfitting
on the limited real-image training set. The learning rate was decayed using a
\texttt{StepLR} scheduler with \texttt{step\_size=3} and \texttt{gamma=0.1}, reducing the
learning rate by a factor of 10 every 3 epochs. All images were resized to
$640 \times 640$ pixels before being fed to the model, and a batch size of 8 was used.

Before committing to the full training run, a series of preliminary training experiments was
conducted on progressively larger subsets of the dataset to characterize the model's learning
behavior and diagnose underfitting and overfitting regimes. These experiments informed the
final hyperparameter choices---most notably the number of training epochs (10) and the
learning rate schedule---by revealing the epoch at which validation loss reached its minimum
before beginning to rise again. The final training run required approximately 174 minutes on
the available hardware, reflecting the computational cost of the two-stage inference
architecture.

\subsection{YOLOv11n}

YOLOv11n was trained using the Ultralytics training framework, which provides a high-level
Python API for model initialization, training loop execution, and checkpointing
\cite{wang2024yolov11}. The nano (n) variant was selected to prioritize inference speed and
deployability on systems without a discrete GPU, while maintaining sufficient model capacity
to learn the relatively constrained visual vocabulary of AST plate images (two object classes
with strong shape priors: circular disks and circular zones).

Model weights were initialized from Ultralytics' official COCO-pretrained checkpoint for
YOLOv11n. Transfer learning from COCO, a large-scale natural image detection benchmark
containing 80 object categories, provides the model with general-purpose object detection
capabilities (feature extraction, non-maximum suppression, anchor-free prediction) that
are then refined for the specific AST domain through fine-tuning. The fine-tuning training
was run for \textbf{200 epochs} on the same dataset split as Faster R-CNN, with images
resized to $640 \times 640$ pixels and a batch size of 8. Built-in Ultralytics data
augmentation was disabled (\texttt{augment=False}) so that the model was trained on
clean, non-augmented images, consistent with the augmentation strategy applied at the
dataset-preparation stage. Training on the full dataset at this configuration required
approximately 200 minutes.

YOLOv11n detects two object classes: \texttt{antibiotic} (the drug-impregnated disk, which
serves as both the OCR target and the in-image calibration reference) and \texttt{disk-zone}
(the surrounding inhibition zone, from which the ZOI diameter is measured). During inference,
the model outputs bounding boxes for all detected objects of both classes, which are then
matched spatially (nearest-disk-to-zone pairing) to associate each zone with its
corresponding antibiotic disk.

\subsection{Hough Circle Transform}

The Hough Circle Transform was applied as a classical baseline using OpenCV's
\texttt{HoughCircles} function, which implements the generalized Hough transform for circle
detection using gradient-based voting. Prior to applying the transform, each image was
converted to grayscale and smoothed with a Gaussian filter using a kernel size of
$17 \times 17$ pixels and $\sigma = 0$ (allowing OpenCV to compute the standard deviation
automatically from the kernel size) to reduce noise that would otherwise produce spurious
gradient responses.

The \texttt{HoughCircles} function accepts four primary tuning parameters: \texttt{dp}
(the inverse resolution ratio of the accumulator to the input image, controlling the
resolution of the parameter-space accumulator), \texttt{minDist} (the minimum allowed
distance between detected circle centers, preventing multiple detections of the same zone),
\texttt{param1} (the upper Canny edge-detection threshold, with the lower threshold set
automatically to half this value), and \texttt{param2} (the accumulator threshold for
center detection, controlling the minimum number of votes required to declare a detected
circle). Optimal values for these four parameters were found by an exhaustive grid search
across the full dataset, minimizing the mean absolute error between detected and ground-truth
circle radii. The final selected values were \texttt{dp=1.5}, \texttt{minDist=150},
\texttt{param1=100}, and \texttt{param2=50}. Importantly, the final reported performance
figures represent the best-case scenario for Hough Transform: the parameter values were
tuned on the full dataset, which likely includes the test set, so generalization performance
may be lower than the reported figures suggest.

\section{ZOI Diameter Measurement and Calibration}

After detecting zones of inhibition as bounding boxes (for the two deep learning models)
or as circles with explicit radius parameters (for the Hough Transform), the raw output must
be converted to a physical ZOI diameter measurement in millimeters before it can be compared
to CLSI or EUCAST breakpoints.

For bounding-box-based detection (Faster R-CNN and YOLOv11n), the pixel-scale diameter of
the detected zone is estimated from the bounding box dimensions using two complementary
approximations that are compared and averaged:
\begin{itemize}
  \item \textbf{Diameter-based:} The average of the bounding box width $w$ and height $h$
        in pixels, $\hat{d}_{\mathrm{diam}} = (w + h)/2$, assuming the zone is approximately
        circular.
  \item \textbf{Area-based:} The equivalent-circle diameter computed from the bounding box
        area, $\hat{d}_{\mathrm{area}} = 2\sqrt{(w \cdot h)/\pi}$, which is more robust
        to mildly elliptical zones where the bounding box aspect ratio deviates from unity.
\end{itemize}

Pixel-to-millimeter conversion is performed using the antibiotic disk as an in-image
calibration reference. Since the physical diameter of a standard CLSI-compliant antibiotic
disk is precisely 6\,mm, and the disk is also detected as a bounding box by the model,
the calibration factor (in millimeters per pixel) can be computed for each image individually
as $k = 6\,\mathrm{mm} / \hat{d}_{\mathrm{disk}}$, where $\hat{d}_{\mathrm{disk}}$ is the
pixel-scale diameter estimate for the corresponding antibiotic disk bounding box. Applying
this per-image calibration factor eliminates the effect of variable camera distance and zoom
level, ensuring that measurements are robust to photographic variation without requiring a
physical ruler in the field of view.

Calibration refinement was an iterative process throughout the project. Initial calibration
using the raw area-based diameter estimate produced a systematic positive bias of approximately
$+4$\,mm (the system consistently over-estimated zone diameters). Analysis of the source of
this bias revealed that bounding boxes predicted by the model systematically overshot the
true zone boundary by a few pixels due to the zone's diffuse, gradient-like edge. By combining
area-based and diameter-based estimates and applying a learned linear correction factor, the
systematic bias was reduced from approximately $+4$\,mm to approximately $+2.1$\,mm,
representing a substantial improvement in measurement accuracy.

\section{OCR Pipeline for Antibiotic Label Recognition}

The OCR pipeline is responsible for reading the alphanumeric antibiotic code from each
detected disk and mapping it to a recognized antibiotic name in the breakpoint database.
The pipeline was designed with a dual-strategy approach to handle the variability in disk
image quality encountered across different laboratory settings and disk brands.

\subsection{Disk Crop and Preprocessing Strategy}

Each antibiotic disk region is first cropped from the full plate image using the bounding box
predicted by the YOLO detection model for the \texttt{antibiotic} class. A small padding
margin is added to the crop boundary to ensure the full disk surface is captured even if the
bounding box slightly underestimates the disk extent.

The dual preprocessing strategy applies OCR in two separate pipelines and selects the result
with the higher confidence score:
\begin{enumerate}
  \item \textbf{Raw image path:} EasyOCR \cite{jaidedai2024easyocr} is applied directly to
        the grayscale disk crop at its original pixel intensity values, without any
        preprocessing transformation.
        This path performs well when the disk image is already high-contrast and the ink
        text is clearly legible without enhancement.
  \item \textbf{CLAHE + Otsu path:} The disk crop is first processed by CLAHE (tile size
        8$\times$8 pixels, contrast limit 2.0) to normalize local illumination variation,
        and then binarized using Otsu's global threshold to produce a high-contrast binary
        image. Morphological dilation or erosion is applied adaptively based on the text
        pixel ratio (the fraction of white pixels in the binarized crop): if the ratio
        falls below 0.05, indicating very thin or faint characters, a $3\times3$ dilation
        kernel is applied to restore stroke connectivity; if the ratio exceeds 0.45,
        indicating over-thickened or merged characters, a $3\times3$ erosion kernel is applied.
        EasyOCR is then applied to this preprocessed binary image.
\end{enumerate}

After obtaining text candidates from both pipelines, the OCR pipeline applies orientation
normalization: EasyOCR is run at \textbf{24 orientations} of the crop (from -180$^\circ$ to +180$^\circ$ in 15$^\circ$ steps) for each pipeline, producing multiple candidate text strings.
This high-resolution rotation strategy ensures that labels can be read regardless of the disk's physical orientation on the agar surface. The candidate with the highest \emph{adjusted} EasyOCR confidence score (factoring in regex pattern matching) across all orientations and pipelines is selected as the final OCR result.

\section{System Deployment and Infrastructure}

The final system is transitioned from a local development environment to a production-equivalent deployment at the \textbf{Chulabhorn Royal Academy (CRA)} hospital server. The production infrastructure is designed for high availability and secure clinical access.

\subsection{Server Environment and Backend Services}

The backend services are hosted on a Linux-based server environment at CRA. The FastAPI application is served using \textbf{Uvicorn}, a lightning-fast ASGI server, running as a managed background process. To ensure system stability, a dedicated \texttt{start.sh} management script handles service orchestration, process monitoring, and automatic recovery. The backend is configured to listen on a high-numbered port (6014) to avoid conflicts with other hospital services.

\subsection{Reverse Proxy and Security}

\textbf{Nginx} is implemented as a reverse proxy and web server to manage incoming traffic to the application. Nginx handles SSL/TLS termination, providing secure HTTPS access to the system. It is configured with specific \texttt{location} blocks to route traffic:
\begin{itemize}
    \item \texttt{/zoneanalyzer/}: Routes to the React frontend (npx serve on port 6015).
    \item \texttt{/zoneanalyzer/api/}: Proxies requests to the FastAPI backend (port 6014).
    \item \texttt{/zoneanalyzer/uploaded_images/}: Serves processed visualization images directly from the server's storage.
\end{itemize}
This configuration enables a single-domain entry point (\texttt{https://paniti-jupyter.cra.ac.th/zoneanalyzer/}) for the entire application, simplifying clinical deployment and firewall management.

\subsection{Frontend Serving}

The React frontend is optimized for production using the Vite build pipeline, which generates a minified, tree-shaken set of static assets. These assets are served using \textbf{npx serve}, providing a lightweight and efficient delivery mechanism for the user interface. The frontend environment is configured with a dedicated \texttt{.env} file pointing to the production API endpoint, ensuring seamless communication across the reverse proxy.

\subsection{Post-Processing and Antibiotic Name Matching}

The raw OCR output is filtered through a regular expression to reject implausible strings.
The regular expression pattern \texttt{\^{}([A-Za-z]+(\textbackslash{}s[A-Za-z]+)\{0,3\})\textbackslash{}s+(\textbackslash{}d+)\$}
enforces the expected format of an antibiotic disk label: one to four alphabetic tokens
(the drug abbreviation) followed by a numeric token (the concentration in micrograms). Strings
that do not match this pattern---which includes OCR hallucinations, partial reads, and pure
numeric strings---are discarded, and the disk is flagged as ``unidentified'' in the results.

Matched antibiotic abbreviations are then looked up in a curated antibiotic name mapping
table that maps standard disk abbreviations (e.g., ``OX'', ``AMX'', ``CIP'') to full antibiotic
names and drug class information. This mapping enables the breakpoint lookup query to retrieve
the correct CLSI or EUCAST threshold for the identified drug and the user-specified organism.

\section{Susceptibility Classification}

The system implements the full CLSI 2025 \cite{clsi2025} and EUCAST 2023 \cite{eucast2023}
breakpoint logic for susceptibility classification. In the Thai clinical laboratory context,
national standards maintained by the Department of Medical Sciences, NIH Thailand
\cite{dmsc_nih} align closely with CLSI guidelines, making CLSI the primary reference for
the target deployment environment. The breakpoint database is stored as a structured SQLite
table with the following schema: antibiotic name, organism group, test standard (CLSI or EUCAST),
disk potency, susceptible threshold diameter ($S_{\min}$ in mm), intermediate range
($I_{\min}$--$I_{\max}$ in mm, if applicable), and resistant threshold diameter ($R_{\max}$ in mm).
This structure directly encodes the published breakpoint tables from both standards.

For each detected zone-disk pair, the system executes the following classification logic:
\begin{enumerate}
  \item Retrieve the breakpoint row matching the OCR-identified antibiotic, the user-specified
        organism group, and the selected standard (CLSI or EUCAST).
  \item Compare the calibrated ZOI diameter $d_{\mathrm{mm}}$ to the retrieved thresholds.
  \item Assign classification: Susceptible (S) if $d_{\mathrm{mm}} \geq S_{\min}$;
        Resistant (R) if $d_{\mathrm{mm}} \leq R_{\max}$; Intermediate (I) if neither
        condition applies (where the Intermediate category exists for the given drug-organism pair).
  \item If no matching breakpoint entry is found for the identified antibiotic and organism
        combination, the classification result is returned as ``Not Available (NA)'' and
        the raw diameter is still reported to allow manual lookup.
\end{enumerate}

Results for all detected disk-zone pairs on the plate are collated and returned to the frontend
in a structured JSON object, which the React dashboard renders as a tabular report listing the
antibiotic name, disk potency, detected zone diameter (in mm), and the S/I/R classification
for each disk. The dashboard also displays an annotated overlay of the original plate image
with bounding boxes drawn around all detected disks and zones, providing a visual confirmation
that the detections are spatially correct. All results are simultaneously persisted to the
SQLite history database for audit and retrospective review.
